\documentclass[10pt,a4paper]{amsart}
\usepackage[margin=19mm]{geometry}
\usepackage{amsmath,amssymb,mathtools}
\usepackage[scr=rsfs,cal=euler]{mathalfa}
\usepackage{xfrac}
\usepackage[unicode,hidelinks,bookmarks=false]{hyperref}
\newcommand{\ie}{i.e.\ }
\newcommand{\cf}{cf.\ }
\newcommand{\resp}{respectively\ }
\newcommand{\Subsection}{\subsection}
\providecommand{\nobreakdash}{\nobreak\hbox{-}}
\setlength{\emergencystretch}{2em}
\multlinegap0pt
\usepackage{bm}
\usepackage{bbm}
\usepackage{dsfont}
\usepackage{xspace}
\usepackage{cancel}
\usepackage{mathtools}
\usepackage{amsthm}
\makeatletter
\@ifundefined{Thm}{\newtheorem{Thm}{Thm}}{}
\@ifundefined{Lem}{\newtheorem{Lem}[Thm]{Lemma}}{}
\makeatother
\title[Crepant resolutions of fiber products and Kummer fibrations ]{Crepant resolutions of fiber products and Kummer fibrations of rational elliptic surfaces with non-reduced fibers}
\author{Dominik Burek, S{\l}awomir Cynk}
\date{Source: arXiv:2609.12522v1 (CC BY 4.0)}
\begin{document}
\maketitle

\noindent\emph{Source and licence.} Crepant resolutions of fiber products and Kummer fibrations of rational elliptic surfaces with non-reduced fibers. Version arXiv:2609.12522v1 (CC BY 4.0), distributed under the Creative Commons Attribution 4.0 International licence, as recorded in the arXiv licence metadata for this exact version and shown on the abstract page https://arxiv.org/abs/2609.12522v1. Full rights notice: licenses/arxiv-2609.12522-CC-BY-4.0.txt.\par\smallskip

\noindent\textit{Editorial scope.} Counted unit: the Lem whose source label is \texttt{None} in the article (source0.tex lines 4725--4727, source label \texttt{None}), delivered together with the article's own complete original proof of it (source0.tex lines 4729--4794, one \texttt{proof} environment of 66 lines). No mathematics is written, repaired or re-derived: statement and proof are the article's own text line for line. Required context retained uncounted: none. Every definition, display and label of those lines is retained unchanged. Omitted from this package and not redistributed: 1--4724, 4728--4728, 4795--5502. Nothing inside a retained excerpt is omitted or altered: every retained body is delivered exactly as the article's lines are written, so no omission receipt of any kind accompanies this package. Citations: the retained excerpts carry no citation command at all, so no bibliography entry of the article is transcribed and no reference is redistributed. Reference adaptation: none was needed and none was made; every reference inside the delivered unit resolves inside it, so no unresolved reference or citation is produced. Printed slips kept literal: no printed word, symbol, hypothesis or number of the counted statement or of its proof was altered. Layout and class: the article's own class is \texttt{amsart} with packages including geometry, latexsym, ifthen, amsmath, enumerate, calc, hyphenat, anyfontsize, moresize, mathtools, bbm, amstext, amsbsy, amsopn; the wrapper is the lane's pinned \texttt{amsart} editorial wrapper at 10pt on A4, which restates the theorem-like environments the retained text needs and adds the article's own macro definitions that the retained text needs (none), with extra wrapper packages (bm, bbm, dsfont, xspace, cancel, mathtools). The delivered numbering is the unit's own. Assets: no retained span contains a figure environment, a graphics-inclusion command, a PDF-page inclusion, a table or a TikZ picture; no image file is copied into this package, the package declares no asset dependency and no image file is redistributed with it. Licence: version arXiv:2609.12522v1 of this article is distributed under the Creative Commons Attribution 4.0 International licence, as recorded in the arXiv licence metadata for this exact version and shown on the abstract page https://arxiv.org/abs/2609.12522v1. A full rights notice is delivered at licenses/arxiv-2609.12522-CC-BY-4.0.txt. Characters: every retained excerpt is pure ASCII, so the delivered engine, XeLaTeX with its default Latin Modern text font, reports no missing character.
\begin{Lem}
Let $X_5$ be the self-fiber product above. After blowing up the ideal $(y-u,x-w)$, the strict transform is smooth in one affine chart, while in the other chart it has a unique ordinary double point. In particular, the local singularity of type $\mathrm{II}\times \mathrm{II}$ admits a crepant resolution obtained by one blow-up followed by a small resolution of a node.
\end{Lem}

\begin{proof}
We analyse the two affine charts of the blow-up of $(y-u,x-w)$.

\smallskip
\noindent
\textbf{First chart.}
Set
\[
y-u=p(x-w).
\]
Then
\[
y=u+p(x-w),
\]
and after subtracting the two defining equations and dividing by $x-w$ we obtain
\[
p(y+u)=x^2+xw+w^2-2t^4.
\]
Substituting $y=u+p(x-w)$ gives
\[
2pu+p^2(x-w)-x^2-xw-w^2+2t^4=0.
\]
Hence the strict transform in this chart is given by
\[
\begin{cases}
u^2-w^3+2t^4w+2t=0,\\
2pu+p^2(x-w)-x^2-xw-w^2+2t^4=0.
\end{cases}
\]
We now examine its singular locus. At the origin
\[
x=w=u=t=p=0
\]
the Jacobian matrix has rank $<2$, and one checks that this is the only singular point in this chart. Moreover, since
\[
\frac{\partial}{\partial t}(u^2-w^3+2t^4w+2t)\bigg|_{0}=2\neq 0
\]
in characteristic $5$, the first equation allows us to solve for $t$ as a formal function of $u$ and $w$. Therefore near the origin the strict transform is analytically a hypersurface in the variables $(x,w,u,p)$, and its quadratic part is
$
2pu-x^2-xw-w^2.
$
This is a non-degenerate quadratic form, so the singularity is an ordinary double point.

\smallskip
\noindent
\textbf{Second chart.}
Set
\[
x-w=q(y-u).
\]
Then
\[
x=w+q(y-u),
\]
and dividing the relation
\[
(y-u)(y+u)=(x-w)(x^2+xw+w^2-2t^4)
\]
by $y-u$ yields
\[
y+u=q(x^2+xw+w^2-2t^4).
\]
Substituting $x=w+q(y-u)$, we obtain the strict transform in this chart. At the origin the Jacobian matrix has rank $2$, hence this chart is smooth.

Therefore, after the blow-up of $(y-u,x-w)$, the strict transform has exactly one singular point, and this singularity is an ordinary double point. A small resolution of this node completes the local crepant resolution.
\end{proof}

\end{document}
